\documentclass[12pt]{article}
\usepackage[margin=1in]{geometry}
\usepackage{amsmath}
\usepackage{amssymb}
\usepackage{xeCJK}
\setCJKmainfont{PingFang SC}

\begin{document}

\section*{Response Ratio / 响应比}

\[
R_X(s)
=
\frac{X'_{2023/24}(s)}
{\overline{X'}_{\mathrm{historical\ strong\ El\ Nino}}(s)}
\]

\noindent
其中：

\begin{itemize}
  \item \(R_X(s)\)：变量 \(X\) 在季节 \(s\) 的响应比。
  \item \(X'_{2023/24}(s)\)：2023/24 事件中，变量 \(X\) 在季节 \(s\) 的异常值。
  \item \(\overline{X'}_{\mathrm{historical\ strong\ El\ Nino}}(s)\)：历史强 El Nino 事件在同一季节 \(s\) 的 composite 异常值。
  \item \(X\)：被比较的变量，例如 Nino-3 SST、u10、SOI-like SLP、GPCP precipitation、CPC OLR。
  \item \(s\)：季节窗口，例如 AMJ、JJA、SON、DJF。
\end{itemize}

\noindent
历史强事件 composite 当前使用：

\[
\mathrm{historical\ strong\ El\ Nino}
=
\{1982/83,\ 1997/98,\ 2015/16\}
\]

\noindent
解释：

\begin{itemize}
  \item \(R_X(s)>1\)：2023/24 在该季节的变量异常强于历史强 El Nino composite。
  \item \(0<R_X(s)<1\)：2023/24 在该季节的变量异常弱于历史强 El Nino composite。
  \item \(R_X(s)<0\)：2023/24 与历史强 El Nino composite 的异常方向相反。
\end{itemize}

\noindent
注意：OLR anomaly 为负值时表示增强对流。因此 CPC OLR 的响应比需要结合符号解释；如果历史强 El Nino 的 OLR 为明显负异常，而 2023/24 的负异常较小，则说明 2023/24 的对流增强偏弱。

\end{document}
